\section*{复数、相量与频率响应}
\noindent 第一编·从二维几何到复数运算\quad\href{course-page:1:6}{1-1}\par
\noindent 第一编·同频振动与相量\quad\href{course-page:1:11}{1-6}\par
\noindent 第十编·电磁波、复数平均与光强\quad\href{course-page:10:85}{10-81}\par
\noindent 第十编·单缝的积分与相量圆弧\quad\href{course-page:10:86}{10-82}\par
\section*{曲率、自然坐标与相对运动}
\noindent 第一编·曲线的切向、法向与曲率\quad\href{course-page:1:76}{1-71}\par
\noindent 第二编·曲率与自然坐标\quad\href{course-page:2:92}{2-90}\par
\noindent 第二编·动坐标与相对运动\quad\href{course-page:2:113}{2-111}\par
\noindent 第三编·梁的弯曲、曲率与端部条件\quad\href{course-page:3:329}{3-326}\par
\noindent 第八编·相对运动、碰撞与自由程\quad\href{course-page:8:15}{8-12}\par
\section*{约束、自由度与广义坐标}
\noindent 第一编·从约束选坐标：抛物线轨道上的小珠\quad\href{course-page:1:81}{1-76}\par
\noindent 第二编·绳杆约束与刚体关联\quad\href{course-page:2:47}{2-45}\par
\noindent 第二编·自由度与坐标选择\quad\href{course-page:2:210}{2-208}\par
\noindent 第三编·能量积分、坐标加速度与约束反力\quad\href{course-page:3:4}{3-1}\par
\noindent 第三编·多自由度、循环坐标与整体反冲\quad\href{course-page:3:21}{3-18}\par
\noindent 第四编·约束运动中的降阶与终止条件\quad\href{course-page:4:75}{4-73}\par
\noindent 第四编·无约束极值与约束极值\quad\href{course-page:4:116}{4-114}\par
\noindent 第五编·局部展开与约束下的稳定性\quad\href{course-page:5:49}{5-47}\par
\section*{虚功、平衡与稳定性}
\noindent 第一编·二次型与平衡点稳定性\quad\href{course-page:1:137}{1-132}\par
\noindent 第二编·平衡点与分界轨线\quad\href{course-page:2:156}{2-154}\par
\noindent 第二编·虚功与静力平衡\quad\href{course-page:2:218}{2-216}\par
\noindent 第三编·约束极值、虚功与二阶变化\quad\href{course-page:3:63}{3-60}\par
\noindent 第三编·平衡构形、二次能量与固有频率\quad\href{course-page:3:135}{3-132}\par
\noindent 第五编·局部展开与约束下的稳定性\quad\href{course-page:5:49}{5-47}\par
\noindent 第八编·相变、界面与外场平衡\quad\href{course-page:8:132}{8-129}\par
\noindent 第八编·相平衡与汽化热\quad\href{course-page:8:142}{8-139}\par
\section*{守恒量与系统边界}
\noindent 第二编·守恒关系与参数平面的允许区域\quad\href{course-page:2:179}{2-177}\par
\noindent 第三编·多自由度、循环坐标与整体反冲\quad\href{course-page:3:21}{3-18}\par
\noindent 第三编·碰撞守恒量、冲量与有效质量\quad\href{course-page:3:180}{3-177}\par
\noindent 第四编·变换中的守恒关系与收敛条件\quad\href{course-page:4:199}{4-197}\par
\noindent 第八编·选择状态变量与守恒的能量\quad\href{course-page:8:132}{8-129}\par
\noindent 第九编·从守恒定律推导火箭方程\quad\href{course-page:9:50}{9-46}\par
\noindent 第九编·四电流与电荷守恒\quad\href{course-page:9:92}{9-88}\par
\noindent 第十编·傅里叶积分中的能量守恒\quad\href{course-page:10:94}{10-90}\par
\section*{相图、势垒与临界运动}
\noindent 第二编·相图的应用\quad\href{course-page:2:147}{2-145}\par
\noindent 第九编·光电效应与表面势垒\quad\href{course-page:9:126}{9-122}\par
\section*{简正模、矩阵与网络谱}
\noindent 第一编·矩阵、线性变换与方程组\quad\href{course-page:1:127}{1-122}\par
\noindent 第一编·非正交坐标与度量矩阵\quad\href{course-page:1:132}{1-127}\par
\noindent 第三编·简正模问题\quad\href{course-page:3:135}{3-132}\par
\noindent 第三编·振型分解、简正坐标与初始条件\quad\href{course-page:3:149}{3-146}\par
\noindent 第七编·有限矩形网络的边界谱与端口电阻\quad\href{course-page:7:82}{7-79}\par
\noindent 第七编·习题2.6：横纵电阻不同的六节点网\quad\href{course-page:7:90}{7-87}\par
\noindent 第十编·约化光线矩阵与复合系统\quad\href{course-page:10:7}{10-3}\par
\noindent 第十编·成像矩阵与横向放大率\quad\href{course-page:10:17}{10-13}\par
\section*{共振、阻尼与能量交换}
\noindent 第一编·受迫阻尼振动：从瞬态到稳态\quad\href{course-page:1:16}{1-11}\par
\noindent 第一编·阻尼模型的完整比较\quad\href{course-page:1:19}{1-14}\par
\noindent 第二编·阻尼及综合应用\quad\href{course-page:2:165}{2-163}\par
\noindent 第三编·拍振、失谐与外力响应\quad\href{course-page:3:159}{3-156}\par
\noindent 第九编·自由原子的吸收共振与反冲\quad\href{course-page:9:154}{9-150}\par
\section*{碰撞、冲量与变质量运动}
\noindent 第二编·连续受力与碰撞散射\quad\href{course-page:2:20}{2-18}\par
\noindent 第三编·碰撞守恒量、冲量与有效质量\quad\href{course-page:3:180}{3-177}\par
\noindent 第三编·多约束、摩擦与连续碰撞\quad\href{course-page:3:196}{3-193}\par
\noindent 第八编·相对运动、碰撞与自由程\quad\href{course-page:8:15}{8-12}\par
\noindent 第八编·碰撞抽样与运动边界的进一步应用\quad\href{course-page:8:27}{8-24}\par
\noindent 第九编·从守恒定律推导火箭方程\quad\href{course-page:9:50}{9-46}\par
\noindent 第九编·火箭的有限点火与燃料分配\quad\href{course-page:9:64}{9-60}\par
\section*{刚体、滚动与进动}
\noindent 第二编·绳杆约束与刚体关联\quad\href{course-page:2:47}{2-45}\par
\noindent 第二编·接触与滚动关联\quad\href{course-page:2:55}{2-53}\par
\noindent 第三编·纯滚动系统的动能与接触条件\quad\href{course-page:3:16}{3-13}\par
\noindent 第三编·滑动摩擦与滚动状态的转换\quad\href{course-page:3:26}{3-23}\par
\section*{微积分、渐近展开与误差}
\noindent 第一编·曲线坐标中的微分算子（选读）\quad\href{course-page:1:101}{1-96}\par
\noindent 第一编·矩阵指数与线性微分方程（选读）\quad\href{course-page:1:155}{1-150}\par
\noindent 第四编·导数、微分与物理过程\quad\href{course-page:4:3}{4-1}\par
\noindent 第四编·泰勒展开与保留阶数\quad\href{course-page:4:12}{4-10}\par
\noindent 第五编·场的微分与曲线坐标\quad\href{course-page:5:13}{5-11}\par
\noindent 第五编·边界误差怎样传入内部区域\quad\href{course-page:5:70}{5-68}\par
\noindent 第六编·局部场方程与微分运算\quad\href{course-page:6:6}{6-3}\par
\noindent 第七编·非均匀泄漏带的幂律与局部近似\quad\href{course-page:7:127}{7-124}\par
\noindent 第九编·角分布积分与四维微分关系\quad\href{course-page:9:27}{9-23}\par
\section*{变分法与最值}
\noindent 第三编·曲线约束与最速下降\quad\href{course-page:3:71}{3-68}\par
\noindent 第四编·二次变分与曲线坐标的几何\quad\href{course-page:4:141}{4-139}\par
\noindent 第十编·费马原理的变分推导\quad\href{course-page:10:13}{10-9}\par
\section*{静电边值、唯一性与镜像}
\noindent 第五编·源、场与电势\quad\href{course-page:5:3}{5-1}\par
\noindent 第五编·边界条件与唯一性\quad\href{course-page:5:46}{5-44}\par
\noindent 第六编·电势、磁矢势与规范\quad\href{course-page:6:7}{6-4}\par
\noindent 第七编·导电与静电的边值类比\quad\href{course-page:7:17}{7-14}\par
\noindent 第七编·唯一性与等势节点的归并\quad\href{course-page:7:47}{7-44}\par
\section*{电介质、极化与磁介质}
\noindent 第五编·极化强度、电位移与束缚电荷\quad\href{course-page:5:132}{5-130}\par
\section*{电磁场能量、动量与应力}
\noindent 第五编·由场能到导体表面的电压力\quad\href{course-page:5:24}{5-22}\par
\noindent 第六编·磁能的形变导数与电磁应力\quad\href{course-page:6:158}{6-155}\par
\section*{电磁感应、磁链与自感}
\noindent 第六编·感应电压的测量路径与滑动接触\quad\href{course-page:6:26}{6-23}\par
\noindent 第六编·磁通、磁链与电感的定义\quad\href{course-page:6:48}{6-45}\par
\section*{粒子轨道与导体运动}
\noindent 第一编·从约束选坐标：抛物线轨道上的小珠\quad\href{course-page:1:81}{1-76}\par
\noindent 第三编·由初态确定轨道的形状与方向\quad\href{course-page:3:356}{3-353}\par
\noindent 第三编·轨道微扰、进动与潮汐\quad\href{course-page:3:377}{3-374}\par
\noindent 第六编·带电粒子的能量与正则动量\quad\href{course-page:6:10}{6-7}\par
\noindent 第六编·磁场边界、转动与受约束轨道\quad\href{course-page:6:160}{6-157}\par
\noindent 第七编·横档与轨道不等阻时的衰减尺度\quad\href{course-page:7:101}{7-98}\par
\noindent 第九编·玻尔氢样原子与相对论圆轨道\quad\href{course-page:9:132}{9-128}\par
\section*{电阻网络的对称性与互易}
\noindent 第一编·立体角、通量与对称性\quad\href{course-page:1:89}{1-84}\par
\noindent 第一编·事件终点的灵敏度与参数优化\quad\href{course-page:1:203}{1-198}\par
\noindent 第六编·互感、互易与耦合强度\quad\href{course-page:6:56}{6-53}\par
\noindent 第六编·几何对称性与互感的分解\quad\href{course-page:6:100}{6-97}\par
\noindent 第七编·对称性之外的桥式网络\quad\href{course-page:7:53}{7-50}\par
\noindent 第七编·互易定理与多端口\quad\href{course-page:7:67}{7-64}\par
\noindent 第七编·习题2.4：电桥失衡后的电流与灵敏度\quad\href{course-page:7:88}{7-85}\par
\noindent 第七编·习题2.5：互易网络接入有限负载\quad\href{course-page:7:89}{7-86}\par
\noindent 第七编·习题2.7：各支路有效电阻的加权和\quad\href{course-page:7:91}{7-88}\par
\noindent 第十编·能流、互易与布儒斯特角\quad\href{course-page:10:132}{10-128}\par
\section*{差分、无限网络与分形}
\noindent 第七编·轮形网络与无限递推\quad\href{course-page:7:75}{7-72}\par
\noindent 第七编·有限梯形与端部条件\quad\href{course-page:7:93}{7-90}\par
\noindent 第七编·习题3.5：两种横档交替的半无限梯形\quad\href{course-page:7:135}{7-132}\par
\noindent 第七编·习题3.6：由有限梯形的测量反求参数\quad\href{course-page:7:136}{7-133}\par
\section*{暂态电路与初值}
\noindent 第五编·边界几何造成的场增强与角模衰减\quad\href{course-page:5:66}{5-64}\par
\noindent 第七编·横档与轨道不等阻时的衰减尺度\quad\href{course-page:7:101}{7-98}\par
\noindent 第七编·串联RLC的二阶暂态\quad\href{course-page:7:156}{7-153}\par
\noindent 第七编·习题4.6：由衰减峰值反求串联电路参数\quad\href{course-page:7:181}{7-178}\par
\noindent 第七编·习题4.7：初始充电选择衰减模态\quad\href{course-page:7:183}{7-180}\par
\section*{分子分布与输运}
\noindent 第一编·连续分布的积分：圆环与圆盘\quad\href{course-page:1:89}{1-84}\par
\noindent 第二编·角分布、立体角与靶面密度\quad\href{course-page:2:12}{2-10}\par
\noindent 第三编·质量输运、动量与能量收支\quad\href{course-page:3:101}{3-98}\par
\noindent 第三编·弯曲波的色散与能量输运\quad\href{course-page:3:333}{3-330}\par
\noindent 第六编·感生电场的空间分布\quad\href{course-page:6:12}{6-9}\par
\noindent 第八编·从状态概率到速度分布\quad\href{course-page:8:4}{8-1}\par
\noindent 第八编·速率、方向与抽样方式\quad\href{course-page:8:9}{8-6}\par
\noindent 第九编·运动光源的角分布\quad\href{course-page:9:16}{9-12}\par
\section*{热力学状态函数与偏导数}
\noindent 第四编·偏导数、全微分与复合函数\quad\href{course-page:4:112}{4-110}\par
\noindent 第八编·泄流、辐射与热力学关系\quad\href{course-page:8:48}{8-45}\par
\noindent 第八编·热力学势与偏导关系\quad\href{course-page:8:57}{8-54}\par
\section*{熵、热机与可用功}
\noindent 第八编·热机、卡诺循环与效率上界\quad\href{course-page:8:90}{8-87}\par
\noindent 第八编·熵的定义与不可逆过程\quad\href{course-page:8:93}{8-90}\par
\section*{相变、相平衡与外场}
\noindent 第八编·相平衡与汽化热\quad\href{course-page:8:142}{8-139}\par
\noindent 第八编·相空间分布与外场中的粒子\quad\href{course-page:8:147}{8-144}\par
\section*{相对论、时空与快度}
\noindent 第三编·固有加速度与相对论推进\quad\href{course-page:3:112}{3-109}\par
\noindent 第九编·共线变换与快度\quad\href{course-page:9:49}{9-45}\par
\noindent 第九编·玻尔氢样原子与相对论圆轨道\quad\href{course-page:9:132}{9-128}\par
\section*{光子、物质波与量子束缚}
\noindent 第五编·极化强度、电位移与束缚电荷\quad\href{course-page:5:132}{5-130}\par
\noindent 第八编·光子气体的压强与辐射\quad\href{course-page:8:55}{8-52}\par
\noindent 第九编·作用量量子化与转折边界\quad\href{course-page:9:133}{9-129}\par
\noindent 第九编·从波长关系到束缚态方程\quad\href{course-page:9:137}{9-133}\par
\section*{几何光学与成像}
\noindent 第十编·成像矩阵与横向放大率\quad\href{course-page:10:17}{10-13}\par
\noindent 第十编·有限孔径、光瞳与可成像的视场\quad\href{course-page:10:24}{10-20}\par
\section*{干涉、衍射与相干}
\noindent 第十编·杨氏干涉与光程差\quad\href{course-page:10:46}{10-42}\par
\noindent 第十编·相干叠加的时间平均\quad\href{course-page:10:57}{10-53}\par
\noindent 第十编·衍射积分、孔径与光栅分辨率\quad\href{course-page:10:85}{10-81}\par
\noindent 第十编·衍射积分中的尺度与归一化\quad\href{course-page:10:93}{10-89}\par
\section*{偏振、晶体与能流}
\noindent 第十编·Jones矢量与偏振椭圆\quad\href{course-page:10:125}{10-121}\par
\noindent 第十编·自然光、部分偏振与检偏\quad\href{course-page:10:127}{10-123}\par
\noindent 第十编·单轴晶体中的普通光与非普通光\quad\href{course-page:10:128}{10-124}\par
\noindent 第十编·全反射相位与倏逝波\quad\href{course-page:10:134}{10-130}\par
\noindent 第十编·晶体中的波矢与能流\quad\href{course-page:10:143}{10-139}\par
\section*{数值计算与参数反演}
\noindent 第一编·简谐振子的数值误差\quad\href{course-page:1:184}{1-179}\par
\noindent 第一编·参数反演中的可识别性\quad\href{course-page:1:199}{1-194}\par
\noindent 第七编·习题3.6：由有限梯形的测量反求参数\quad\href{course-page:7:136}{7-133}\par
\noindent 第七编·习题4.6：由衰减峰值反求串联电路参数\quad\href{course-page:7:181}{7-178}\par
